\documentclass[11pt,a4paper]{article}
\usepackage[margin=0.78in]{geometry}
\usepackage[T1]{fontenc}
\usepackage{lmodern,amsmath,amssymb,booktabs,tabularx,xcolor}
\usepackage[hidelinks]{hyperref}
\usepackage{xurl}
\usepackage{enumitem}
\setlist{nosep,leftmargin=*}
\setlength{\parindent}{0pt}
\setlength{\parskip}{4pt}
\newif\ifguidance
\guidancefalse
\newcommand{\help}[1]{\ifguidance{\small\color{blue!55!black}\textit{Writing guide: #1}\par}\fi}
\newcommand{\fillin}[1]{\textnormal{[\textit{#1}]}}
% Team facts used by every document in writeup/, report/ and slides/. Edit here once.
\newcommand{\StudentA}{Sumedh Navuda}
\newcommand{\StudentB}{Aldino}
\newcommand{\Authors}{\StudentA{} and \StudentB{}}
\newcommand{\Contact}{sumedhnavuda@outlook.com}
\newcommand{\RollA}{[roll no.]}   % <- replace with Sumedh's roll number
\newcommand{\RollB}{[roll no.]}   % <- replace with Aldino's roll number
\newcommand{\RepoURL}{https://github.com/sumedhnvda/ict5143-p8-fairness-audits}
\newcommand{\TeamID}{[team ID]}   % <- replace with the team ID given by the instructor
\newcommand{\SubmitTag}{submission-v1}

\newcommand{\ProjectTitle}{When Does Pooling Make Intersectional Fairness Audits Falsely Reassuring?}
\begin{document}\sloppy
\begin{center}
{\Large\bfseries \ProjectTitle}\\[4pt]
ICT 5143: Algorithms Fairness and Graphical Models Laboratory\\
Problem 8 (Reliable intersectional fairness audits); Team \TeamID\\
\StudentA{} (\RollA); \StudentB{} (\RollB)\\
{\small Repository: \url{\RepoURL}; contact: \texttt{\Contact}\\
Submission tag: \texttt{\SubmitTag}; Date: 6 November 2026}
\end{center}
\begin{abstract}
We ask when Bayesian pooling makes an intersectional equal-opportunity audit report ``gap $\le0.10$'' for an
outreach rule that is unfair. A discrete Bayesian network ranks ACSPublicCoverage adults (California 2018); audits
estimate race$\times$sex$\times$disability TPRs. Pooled audits hide isolated small-group deficits, and all 116 California
and 27 Texas validation-time tolerance claims failed on test. Calibrating the audit on the target population's group
sizes and certifying on held-out households raised the audit score on fresh Texas simulations by $+0.279$
(95\% CI $[0.235,0.322]$). Limitation: on California it only ties the best pooled audit.
\end{abstract}

\section{Question, context and contribution}
An agency with capacity for 30\% of a list sends \emph{hypothetical} information outreach ($D=1$) to people ranked
most likely to lack public coverage ($Y=1$: no public coverage in the ACS record; this is neither uninsurance nor
programme eligibility). The people affected are those who lack coverage but are never contacted, often in small
intersections such as disabled Black women. The auditor must say whether reach is equal across groups.
\textbf{Research question:} when does hierarchical pooling stabilise intersectional TPR estimates, and when does it
mask a harmful small-group disparity, so that a tolerance claim $\Delta\text{EOpp}\le\epsilon$ is false?\\
\textbf{Contribution:} pooled audit criteria used to choose a group-threshold rule on validation data make tolerance
claims that fail on test in more than half of cases (H1; refuted if at most half fail on Texas). Design-matched
calibration with held-out certification improves the audit score $J$ by at least 0.10 over every established audit
(H2; refuted if not).\\
\textbf{Closest work:} hierarchical Bayesian intersectional fairness~\cite{foulds2020bayes} and structured regression
for small subgroups~\cite{herlihy2024}. Our \texttt{struct} audit is a restricted version of the former, and we do not
claim an exact reproduction. The contamination prior is classical~\cite{berger1986}, so we claim an evaluation and an
audit protocol, not a new prior. The 19-paper comparison is in \texttt{writeup/literature\_review.pdf}.

\section{Data, graphical model and fairness method}
\textbf{Data and estimand:} ACS PUMS 2018 1-year person file via Folktables ACSPublicCoverage~\cite{ding2021}
(AGEP\,$<$\,65, PINCP\,$\le$\,30{,}000): 138{,}554 Californians in 85{,}182 households, prevalence $P(Y{=}1)=0.631$.
Label $Y=1-\mathbf{1}[\text{PUBCOV}=1]$. Protected attributes are the recorded race (9 codes), sex and disability.
Their 36 intersections are the audit groups. Splits are 60/20/20 \emph{by household} (SERIALNO) over 5 seeds. The
estimand is the unweighted benchmark sample; survey weights are a robustness check. Texas 2018 (98{,}928) is the
replication population.\\
\textbf{Graph and inference:} an augmented naive-Bayes network (associational, not causal). $Y$ is a parent of all 18
features, with ten feature edges (e.g.\ DIS$\mid$AGEP, PINCP$\mid$ESR) and Dirichlet $\alpha=1$ CPTs. The audit is a
second graphical model over group counts: $k_g$ selected positives out of $n_g$ positives.
\begin{gather*}
p(y,\mathbf{x})=p(y)\prod_{j}p(x_j\mid\operatorname{pa}(x_j)),\qquad s(\mathbf{x})=P(Y{=}1\mid\mathbf{x})\ \text{(variable elimination)},\\
k_g\mid\theta_g\sim\text{Bin}(n_g,\theta_g),\qquad \theta_g\sim\pi\,\text{Beta}(a,b)+(1-\pi)\,\text{Beta}\big(\mu c,(1-\mu)c\big),\quad \pi\le\pi_{\max}.
\end{gather*}
Comparators replace the prior on $\theta_g$: none (raw), independent Beta(1,1), complete pooling, exchangeable
Beta$(a,b)$ by empirical Bayes or a full-Bayes grid, and \texttt{struct}:
$\operatorname{logit}\theta_g=\beta_0+\beta_{\text{race}}+\beta_{\text{sex}}+\beta_{\text{dis}}+u_g$,
$u_g\sim N(0,\tau^2)$, fitted by quadrature.\\
\textbf{Method and fairness:} the primary criterion is intersectional equal opportunity,
$\Delta\text{EOpp}=\max_g\text{TPR}_g-\min_g\text{TPR}_g$, with $\epsilon=0.10$. It fits outreach because the harm
is missing people who need it. Diagnostics are $\Delta$DP, $\Delta$EOdds, selection-rate ratio, worst-group recall and
group calibration. Utility is recall at 30\% capacity, with at most 0.02 loss allowed. A rule family $R(\lambda)$
blends a global and group-specific thresholds, so $A$ is used at prediction time for $\lambda>0$ (the scores also use
$A$). An audit certifies when its plug-in or upper-$q$ posterior gap is $\le\epsilon$. \textbf{Proposed:} choose the
pooling cap $\pi_{\max}$, slab and rule to maximise $J=\text{ACC}-\overline{\text{FTS}}$ on semi-synthetic worlds
built on the \emph{target} state's validation group sizes. Then choose $\lambda$ on half of the validation households
and certify on the other half.

\section{Evaluation protocol}
Baselines are logistic regression and gradient boosting (predictive), the BN with no uncertainty-aware selection
(graphical), and the raw, independent, exchangeable, complete and \texttt{struct} audits (fairness). All audits see
the same counts, and all rules share the capacity matched on validation. Discretisation edges and CPTs are fitted on
training; $\lambda$ and thresholds are chosen on validation; test is used only to evaluate frozen rules. Audit
procedures are ranked by $J$ on simulated worlds with known truth: homogeneous (fair) worlds and three unfair worlds
with a 0.20 TPR deficit in an isolated cell, a margin-balanced interaction or a disability main effect, at full and
half support. Fresh simulation seeds are used for evaluation. Five household splits; household-cluster Poisson
bootstrap (B\,=\,1000) for real paired gaps, with Holm correction over the two preregistered comparisons. $\Delta J$
uses a paired bootstrap over simulated worlds (B\,=\,2000, 200 worlds per cell). Exact VE is checked against pgmpy
VE, clique-tree propagation and enumeration (tolerance $10^{-9}$).

\section{Results and interpretation}
\begin{center}\small
\begin{tabularx}{\linewidth}{Xcccc}
\toprule
Audit procedure (fresh Texas worlds, 200 per cell) & ACC (fair, $\uparrow$) & mean FTS ($\downarrow$) & $J$ ($\uparrow$) & 90\% gap coverage\\
\midrule
Closest baseline: full-Bayes beta-binomial, plug-in & 1.000 & 0.615 & 0.385 & 0.174\\
Hierarchical EB, upper 0.9 & 0.985 & 0.625 & 0.360 & 0.013\\
\texttt{struct} (Foulds-style), plug-in & 0.108 & 0.073 & 0.034 & 0.037\\
Ablation: no design matching (\S9 calibration) & 0.218 & 0.028 & 0.189 & 0.506\\
\textbf{Proposed: design-matched contamination audit} & 0.832 & 0.168 & \textbf{0.664} & 0.615\\
\midrule
Real CA+TX, 10 splits & \multicolumn{2}{c}{claims / false claims} & test recall & test $\Delta$EOpp\\
Ablation: same-data certification & \multicolumn{2}{c}{10 / 10} & 0.365 & 0.302\\
\textbf{Proposed: held-out certification} & \multicolumn{2}{c}{0 / 0} & 0.365 & 0.309\\
\bottomrule
\end{tabularx}
\end{center}
{\small \textbf{Table note:} ACC is the share of fair worlds certified, and FTS the share of unfair worlds certified
(false tolerance satisfaction), averaged over three deficit types and two support levels at $\epsilon=0.10$.
Paired $\Delta J$ (proposed minus baseline) is $+0.279$ [0.235, 0.322] against full Bayes and $+0.475$ [0.422, 0.525]
against the ablation. On the California pilot design it is $-0.004$ [$-0.057$, 0.046], a tie. Predictive quality is
shared by all audits: the BN has AUROC $0.729\pm0.004$, log loss 0.577 and Brier 0.196 (GBT: 0.768, 0.541, 0.182).}

\textbf{Main finding:} pooling helps when deficits are carried by large groups or main effects, and it masks them
when they are isolated. On real data the BN rule's disparity is structural, and every audit gives
$P(\Delta\text{EOpp}>0.10)=1$. The gap grows from 0.162 by sex to 0.746 over the 21 supported intersections (35
observed). On synthetic worlds, EB with the upper-0.9 rule falsely certified 98\% of worlds with one 20-positive group
0.20 below the rest. In selection, all 116 California and 27/27 Texas pooled validation claims failed on test (test
gap 0.110--0.888), so \textbf{H1 is supported}. The proposed audit raises $J$ on Texas by more than 0.10 over every
established audit and over the ablation (\textbf{H2a, H2b supported}). Held-out certification withheld all 10 false
claims at unchanged recall (\textbf{H2c}).\\
\textbf{Mechanism and failure:} removing design matching drops fair-world acceptance from 0.83 on invented sizes to
0.22 on Texas. 90\% of false alarms come from extreme groups with fewer than 10 positives, so the calibration must
see the real group sizes. Failures: on California the method trades lower FTS (0.348 vs 0.666) for lower ACC (0.678),
which leaves $J$ unchanged. No rule in the family is fair on test (minimum test gap 0.227). The preregistered
real-gap comparison is inconclusive ($-0.052$ [$-0.194$, 0.051], Holm $p=0.736$). The \texttt{struct} intervals
under-cover (0.665 at nominal 0.90; oracle 0.905). Groups with fewer than 30 positives are exploratory.

\section{Limitations, ethics and conclusion}
The ACS facts are the counts, gaps and failed claims. The $J$ gains are simulation evidence from one generator family
on real group sizes. Race and disability are survey constructs, and the label is not need, so a model trained on it
can misdirect resources. Zero held-out claims shows abstention, not sensitivity to rules that really are fair. No
stakeholders set $\epsilon$ or the capacity. \textbf{Conclusion:} a pooled ``gap $\le\epsilon$'' should not be
trusted for sparse intersections unless the audit was calibrated on that design and certified on households not used
to choose the rule. Next experiment: a generator family not used in calibration, and a constructed fair rule to
measure held-out sensitivity.

\section*{Reproducibility and individual contributions}
\textbf{Reproduce:} \texttt{jupyter nbconvert -{}-to notebook -{}-execute -{}-inplace} \texttt{Problem8\_\allowbreak Intersectional\_\allowbreak Fairness.ipynb}
(about 30 min on 1 CPU core; Python 3.12, \nolinkurl{requirements-lock.txt}), then \texttt{python writeup/make\_tables.py}
for the $\Delta J$ intervals. Config: notebook \S0 constants, mirrored in \nolinkurl{configs/frozen_settings.yaml}.\\
\textbf{\StudentA:} question and preregistration, 19-paper review, data and household splits, BN and VE/clique-tree
verification, Experiment 1, Texas H1 test. Cross-checked the audit posteriors, Holm and $J$ intervals.\\
\textbf{\StudentB:} audit models and coverage checks, simulations and calibration, threshold selection, bootstrap,
the design-matched/held-out procedure (H2). Cross-checked split disjointness, label reversal and inference agreement.\\
\textbf{Evidence:} \texttt{docs/requirements\_matrix.md}, \texttt{docs/contributions.md}, \texttt{results/}.
\begin{thebibliography}{9}\small
\bibitem{foulds2020bayes} J.~R. Foulds, R.~Islam, K.~N. Keya, S.~Pan. Bayesian modeling of intersectional fairness: the variance of bias. \emph{Proc.\ SIAM SDM}, 424--432, 2020. doi:10.1137/1.9781611976236.48.
\bibitem{herlihy2024} C.~Herlihy, K.~Truong, A.~Chouldechova, M.~Dud\'ik. A structured regression approach for evaluating model performance across intersectional subgroups. \emph{Proc.\ ACM FAccT}, 2024. arXiv:2401.14893.
\bibitem{berger1986} J.~O. Berger, L.~M. Berliner. Robust Bayes and empirical Bayes analysis with $\epsilon$-contaminated priors. \emph{Ann.\ Statist.} 14(2):461--486, 1986.
\bibitem{ding2021} F.~Ding, M.~Hardt, J.~Miller, L.~Schmidt. Retiring Adult: new datasets for fair machine learning. \emph{NeurIPS 34}, 6478--6490, 2021. Data: U.S.\ Census Bureau ACS PUMS 2018 1-year, via \url{https://github.com/socialfoundations/folktables}, accessed October 2026.
\end{thebibliography}
\end{document}
